\section*{Introduction}

L'étude de la pensée politique d'Ambroise de Milan doit nécessairement se faire par une première approche théorique à travers une analyse des fondements intellectuels et théologiques des écrits ambrosiens. Ce premier chapitre a pour objectif d'examiner l'ensemble de l'œuvre de l'évêque pour y déceler les éléments qui constituent sa théorie du pouvoir et du gouvernement romain. L'ensemble de sa conception sur un pouvoir chrétien, et donc l'ensemble de ses relations avec la cour impériale, repose sur un questionnement concernant l'organisation de la société et le rôle des citoyens en plus de l'agencement des détenteurs du pouvoir, ne se fermant pas derrière un seul concept politique. Il se montre ainsi capable d'identifier les aspects positifs de diverses formes de gouvernement, parlant aussi bien des systèmes de partage de la république des grues que de l'absolutisme des abeilles et du monde romain.

\bigskip

Cette diversité au sein de sa pensée est d'ailleurs mise en évidence par l'historiographie, avec entre autres Marc Reydellet dans ses travaux sur les souverains face à la Bible qui souligne l'ambiguïté ambrosienne sur les fondements du pouvoir, l'évêque loue les principes républicains d'une alternance du pouvoir, mais accepte sans problème le système monarchique de son temps\autocite{reydellet_bible}. Ce postulat qui peut sembler contradictoire trouve en réalité son sens dans son interprétation de l'épître aux Romains qui démontre que tout pouvoir vient de Dieu, et donc que c'est l'utilisation humaine qui est mauvaise. Si le fait même de posséder le pouvoir n'entre pas en contradiction avec la foi chrétienne, alors le bon gouvernement peut trouver sa place derrière le système monarchique, il suffit de le guider correctement. À partir de ce cadre, Ambroise réfléchit aux atouts de l'utopie de la \latin{res publica} pour les transférer dans un souverain unique tout aussi dévoué à l'intérêt général.

\bigskip

Cette réflexion, bien que théorique, se met en place autour des nombreux rôles joués par Ambroise tout au long de son épiscopat. En effet, notre analyse de sa pensée doit se faire sur des textes qui sont pour la plupart voués à être prononcés devant un auditoire, mais celui-ci peut être très varié, ce qui implique une grande adaptation d'Ambroise qu'il soit évêque devant son clergé, prédicateur devant le peuple ou conseiller politique devant l'empereur. En fonction de sa posture, son discours évolue et il est nécessaire de recouper ces différentes approches pour en tirer l'ensemble de sa théorisation du bon gouvernement. Comme ses textes se doivent d'être utiles, sa pensée se fonde davantage sur une morale chrétienne de la gestion de la société que sur l'élaboration de nouvelles règles constitutionnelles. Sa façon de théoriser le gouvernement s'intègre dans une vision plus large de l'exercice du pouvoir : il cherche à établir une proximité et une confiance entre les détenteurs du pouvoir, peu importe le régime, et le peuple, en considérant que cette cohésion est la seule véritable condition à la bonne tenue d'un gouvernement et d'une société juste.

\bigskip

Cette théorie politique trouve son fondement dans de multiples observations des sociétés animales, de sa lecture des textes bibliques et de son analyse des régimes politiques humains, républicains ou monarchiques. Pour pleinement réussir à comprendre la pensée ambrosienne, il est nécessaire d'observer avec attention plusieurs de ses œuvres. C'est par exemple le cas du \latin{De Officiis} qui fait office de longue réflexion sur les devoirs des citoyens au sein de la vie en communauté et fournit plusieurs éléments importants de sa théorie politique. C'est également le cas de son \latin{Exameron} et plus particulièrement son commentaire sur le cinquième jour de la Création qui apporte l'essentiel de sa pensée sur le modèle républicain et l'idéal de partage et de mise en commun du pouvoir et des tâches pour pousser tous les membres d'une société dans la même direction. Enfin, ses oraisons funèbres en l'honneur de Valentinien II et Théodose nous permettent de saisir les attentes concrètes d'Ambroise concernant le gouvernement romain, en mettant en avant les réussites des actions des deux empereurs défunts.

\bigskip

Ainsi, ce chapitre cherchera dans un premier temps à analyser ce qui constitue l'idéal républicain d'Ambroise de Milan pour comprendre les éléments qu'il cherche à introduire dans sa théorisation du gouvernement romain. Nous nous concentrerons alors sur son observation du monde naturel pour étudier la façon dont il conçoit un régime politique fondé sur la mise en commun et le partage des responsabilités. Dans une deuxième partie, nous continuerons l'analyse de sa théorie du gouvernement romain à travers son pragmatisme d'un modèle monarchique qu'il accepte mais non sans apporter des changements et limites importants, porté par la morale chrétienne, afin de rapprocher le système impérial romain de sa théorie utopique.

\bigskip
